\chapter{毕业设计相关关键技术}

基于树莓派的人脸识别考勤系统涉及端侧计算、计算机视觉、Web服务、本地数据存储和远程组网等多类技术。树莓派提供低成本边缘部署环境，摄像头采集模块提供实时图像输入，OpenCV YuNet与SFace分别承担人脸检测和特征提取任务，OpenVINO用于部署活体检测模型，FastAPI、Jinja2和Bootstrap支撑页面与接口服务，SQLite用于保存本地业务数据，Tailscale用于构建远程管理访问通道。这些技术共同构成系统实现的基础。

\section{树莓派端侧部署与摄像头采集}

树莓派是一类基于ARM架构的单板计算机，能够运行Linux操作系统，并提供USB、以太网、无线网络、GPIO和本地存储等基础能力。与普通服务器相比，树莓派体积小、功耗低、部署成本低，适合放置在教室门口、实验室入口、办公室等空间中承担轻量化边缘计算任务。端侧部署的核心思想是把数据采集、初步处理和部分业务决策放在靠近数据源的位置完成，减少对中心服务器和公网链路的依赖。对于人脸识别考勤场景而言，摄像头画面包含较强的个人隐私属性，如果所有视频帧都上传到远程服务器处理，会增加带宽压力和隐私风险；而在树莓派端完成采集、识别和写库，则可以让系统在小规模场景中保持更强的自治能力。

摄像头采集是端侧视觉系统的基础。Raspberry Pi官方文档提供了摄像头软件栈和图像采集相关说明\cite{raspberrypi_camera}。在工程实现中，USB摄像头通常以UVC设备形式被Linux识别，应用程序可以通过OpenCV的\texttt{VideoCapture}接口打开摄像头、设置分辨率和帧率，并循环读取图像帧。读取到的原始帧一般还需要经过尺寸限制、方向旋转、镜像处理、颜色空间转换和图像编码等步骤，才能同时满足模型推理和浏览器显示的需要。

端侧部署对系统提出了两类约束。第一类是算力约束。树莓派CPU资源有限，人脸检测、特征提取和活体检测都不能无限制占用计算资源，否则会导致视频流卡顿、页面响应变慢，甚至影响服务稳定性。第二类是长期运行约束。考勤终端往往需要连续运行，系统必须处理摄像头打开失败、读帧失败、网络临时不可用、远程管理端离线等情况。基于上述特点，本文系统使用普通USB摄像头作为图像输入设备，后端通过OpenCV读取视频帧，并将视频流以MJPEG形式输出到本地终端页面；在实现上采用采集线程与识别线程解耦的方式，采集线程持续读帧和推流，识别线程只处理共享缓冲区中的最新帧，从而避免识别任务形成帧积压。

\section{OpenCV YuNet人脸检测模型}

人脸检测是人脸识别链路的前置环节，其任务是在图像中定位人脸区域，并给出后续对齐和特征提取所需的关键点信息。一个完整的人脸识别系统通常不能直接对整张图像提取身份特征，因为图像中可能包含背景、身体、多人、遮挡物和不同尺度的人脸。人脸检测模型首先从图像中筛选出可能的人脸区域，再输出边界框、关键点和置信度，为后续的人脸对齐、质量判断和身份识别提供输入。

YuNet是面向轻量化人脸检测任务设计的模型，论文中将其定位为小型毫秒级人脸检测器\cite{wu2023yunet}。与体积较大的通用目标检测模型相比，YuNet更关注人脸这一特定目标，在模型规模、推理速度和检测精度之间取得平衡，适合部署在CPU设备或边缘设备上。YuNet不仅输出人脸边界框，还输出人脸关键点。关键点通常包括左右眼、鼻尖和嘴角等位置，这些点可用于判断人脸姿态、执行人脸对齐以及估计图像质量。

OpenCV官方文档将YuNet与后续SFace识别流程组织成DNN人脸检测与识别示例\cite{opencv_dnn_face}。在OpenCV中，YuNet检测器通过\texttt{FaceDetectorYN}接口创建和调用。OpenCV文档说明该接口用于创建、配置和执行人脸检测器，并可设置输入尺寸、置信度阈值、NMS阈值和最大候选数\cite{opencv_facedetectoryn}。其中，置信度阈值用于过滤低可信候选框，NMS阈值用于消除高度重叠的重复检测框，最大候选数用于限制模型输出规模。这些参数对端侧实时系统非常重要，因为阈值过低会带来误检和额外计算，阈值过高则可能漏检真实人脸。

设第$i$个候选人脸检测结果为：

\begin{equation}
    D_i = (x_i, y_i, w_i, h_i, l_{i1}, l_{i2}, \dots, l_{i10}, s_i)
\end{equation}

其中，$(x_i, y_i)$表示边界框左上角坐标，$w_i$和$h_i$分别表示边界框宽度和高度，$l_{i1}$至$l_{i10}$表示5个关键点的二维坐标，$s_i$表示检测置信度。检测结果进入业务流程后，还需要结合考勤场景进行二次判断。例如，当画面中没有人脸时，系统应提示用户进入画面；当检测到多张人脸时，系统应暂停签到，避免身份归属不清；当人脸区域过小时，系统应提示用户靠近摄像头。本文系统采用YuNet作为最终人脸检测模型，用于终端识别中的人脸定位，也为后续SFace特征提取、注册照质量校验和姿态估计提供基础输入。

\section{OpenCV SFace特征提取模型}

人脸识别的关键是将人脸图像转化为可比较的特征表示。早期方法常使用人工设计特征描述人脸纹理或局部结构，但在光照变化、姿态变化、年龄变化和表情变化下鲁棒性有限。深度学习方法通过神经网络从大量人脸样本中学习特征映射，将一张人脸图像表示为固定维度的向量。理想情况下，同一身份的不同图像在特征空间中距离较近，不同身份的图像距离较远。这样，系统就可以通过向量相似度完成身份匹配。

SFace是一种基于Sigmoid约束超球面损失的人脸识别方法，能够提升人脸特征在相似度空间中的可分性\cite{zhong2021sface}。该类方法通常会将特征向量约束在归一化空间中，使余弦相似度可以直接反映两个样本的身份接近程度。与直接使用分类概率不同，人脸考勤系统更适合使用特征向量匹配方式，因为系统中的人员集合会随着成员注册和审核不断变化，不适合每增加一个人就重新训练分类器。特征向量方式只需要将审核通过的人脸向量保存到数据库中，识别时把当前人脸向量与库中向量逐一比较即可。

OpenCV提供的\texttt{FaceRecognizerSF}接口支持人脸对齐、特征提取和特征匹配等操作\cite{opencv_facerecognizersf}。人脸对齐是特征提取前的重要步骤。由于用户面对摄像头时可能存在轻微旋转、平移或尺度差异，直接对原始人脸框提取特征会增加不必要的变化。对齐过程利用检测模型输出的关键点，将人脸裁剪和归一化到模型期望的姿态和尺度，从而提高特征稳定性。完成对齐后，SFace模型输出人脸特征向量，系统再对向量进行归一化处理，以便进行余弦相似度计算。

设原始特征向量为$f$，归一化后的特征向量为$\hat{f}$：

\begin{equation}
    \hat{f} = \frac{f}{\|f\|_2}
\end{equation}

其中，$\|f\|_2$为向量的二范数。归一化后，两个特征向量之间的余弦相似度可以通过点积计算：

\begin{equation}
    \text{sim}(\hat{f}_1,\hat{f}_2) = \hat{f}_1 \cdot \hat{f}_2
\end{equation}

相似度越高，表示两张人脸越可能属于同一人。在本文系统中，SFace负责把YuNet检测并对齐后的人脸转化为特征向量。终端识别时，系统将当前帧人脸特征与数据库中所有审核通过的人脸特征进行比较，选择相似度最高的记录作为候选身份；当相似度不低于配置阈值0.363时，系统判定身份匹配通过；当候选结果过于接近时，系统返回歧义状态，避免误签到。

\section{OpenVINO与活体检测技术}

活体检测是人脸识别系统安全性的重要组成部分。单纯的人脸检测和特征匹配只能回答“画面中的人脸是否像某个人”，不能回答“画面中的人脸是否来自现场真人”。在考勤系统中，常见攻击包括纸质照片、手机屏幕翻拍、平板播放照片或视频等。如果系统只依据相似度写入考勤记录，就可能被静态照片或屏幕重放欺骗。因此，活体检测需要在身份匹配之外增加一层真实性判断。

活体检测方法可以分为主动式和被动式两类。主动式活体检测要求用户完成眨眼、摇头、张嘴、转头等动作，然后由系统判断动作是否符合要求。这类方法交互明确，适合注册采集等低频场景，但会增加用户操作时间。被动式活体检测不要求用户额外动作，而是从单帧或多帧图像中提取纹理、反光、边缘、频域或模型特征，判断当前画面是否可能来自真实人脸。相关综述指出，深度学习活体检测需要关注跨设备、跨光照和跨攻击介质的泛化能力\cite{yu2023fas}。在实际工程中，被动式方法体验更自然，但容易受到摄像头质量和现场光照影响，因此通常需要结合多帧统计或其他风险信号。

OpenVINO是Intel提供的深度学习模型优化和推理工具套件，支持在CPU等硬件上部署模型推理任务\cite{openvino}。它通常将模型转换或加载为适合推理引擎执行的表示，并通过设备插件在CPU等硬件上运行。对于树莓派这类资源有限的设备，推理框架的价值在于减少模型部署成本，使应用程序可以以较统一的方式加载模型、准备输入张量、执行推理并读取输出结果。

本文使用OpenVINO加载Open Model Zoo中的anti-spoof-mn3模型。该模型文档说明其用于区分真人人脸和伪造人脸，输入为裁剪后的人脸图像，输出为真人与伪造两类概率\cite{openvino_anti_spoof}。系统首先根据YuNet输出的人脸框裁剪人脸区域，然后将图像缩放到模型输入尺寸，并转换为模型需要的张量格式。模型输出经过归一化后得到真人概率$S_{\text{real}}$和伪造概率$S_{\text{spoof}}$，满足：

\begin{equation}
    S_{\text{real}} + S_{\text{spoof}} = 1
\end{equation}

在考勤终端中，单帧活体结果容易受运动模糊、光照突变和裁剪偏差影响。因此，本文系统对单帧结果采用三区间判断策略：当$S_{\text{real}} \geq 0.55$时，认为当前帧活体可信；当$0.30 \leq S_{\text{real}} < 0.55$时，认为当前帧处于灰区，需要结合多帧结果继续判断；当$S_{\text{real}} < 0.30$时，记为低分事件。系统还维护长度为2秒的滑动窗口，对窗口内样本进行统计。设窗口内共有$M$个样本，其真人概率为$S_1,S_2,\dots,S_M$，则平均真人概率为：

\begin{equation}
    \bar{S}_{\text{real}} = \frac{1}{M}\sum_{i=1}^{M}S_i
\end{equation}

当窗口样本数不足或分数不稳定时，系统提示用户继续保持正对摄像头；当低分次数达到配置上限，或屏幕重放风险超过阈值时，系统拒绝签到。这样既保留了被动活体检测的低交互优势，又通过多帧融合降低单帧误判风险。

\section{FastAPI、Jinja2与Bootstrap}

Web服务框架负责将后端业务能力以页面和接口形式暴露给用户。本文系统需要同时提供本地终端页面、用户门户、动作挑战接口、管理员API、图片访问接口和视频流接口，因此后端框架需要具备路由清晰、接口组织方便、请求响应类型丰富、易于部署等特点。FastAPI是基于Python类型提示构建的Web框架，官方文档强调其支持自动生成交互式API文档、请求数据校验和异步能力\cite{fastapi}。FastAPI建立在ASGI生态之上，可以通过Uvicorn运行，适合处理普通JSON接口、HTML页面响应、文件响应和流式响应等不同类型的请求。

FastAPI的一个重要特点是使用“路径操作函数”组织接口。开发者可以通过装饰器定义HTTP方法和URL路径，再在函数参数中声明查询参数、路径参数、请求对象或数据模型。配合Pydantic，FastAPI能够对请求和响应数据进行结构化校验，这有助于减少接口输入格式错误。在本文系统中，FastAPI用于组织三类入口：本地终端入口、普通用户门户入口和管理员API入口。终端入口提供根页面、MJPEG视频流和识别状态接口；用户门户入口提供注册、登录、人脸提交和动作挑战接口；管理员入口提供设备状态、成员管理、人脸审核和考勤查询接口。

Jinja2是Python生态常用模板引擎，官方文档说明其支持模板继承、变量输出和控制结构\cite{jinja}。模板引擎的作用是把后端数据与HTML结构结合起来，生成浏览器可直接渲染的页面。Jinja2支持公共布局抽取、页面块覆盖、条件判断、循环输出和过滤器等能力，适合中小型管理页面和业务表单页面。与单页应用框架相比，服务端模板不需要额外前端构建流程，部署更简单，也更适合树莓派端这种轻量场景。

Bootstrap提供响应式布局、表单、按钮和组件样式，其官方文档将其定位为前端工具包\cite{bootstrap}。它通过栅格系统、表单控件、按钮样式、卡片和工具类降低基础页面开发成本。本文系统使用Jinja2生成终端识别页面、用户注册页面、登录页面、用户中心页面和人脸采集页面，使用Bootstrap构建页面基础布局，使终端页面和用户门户在视觉和交互上保持一致。对于动作挑战和终端状态刷新等动态功能，系统使用原生JavaScript调用后端接口，不引入额外前端构建工具，从而降低部署复杂度。

\section{SQLite数据库}

SQLite是一种嵌入式关系型数据库，官方文档说明其不需要独立服务器进程，数据库以普通磁盘文件形式保存\cite{sqlite}。与MySQL、PostgreSQL等客户端/服务器模式数据库不同，SQLite直接作为应用程序的一部分运行，应用通过本地文件完成数据读写。它支持SQL语句、事务、索引、外键和多种日志模式，适合数据规模不大、部署环境简单、需要本地持久化的应用场景。

SQLite的优点主要体现在部署和维护成本低。对于树莓派端的小规模考勤系统而言，使用独立数据库服务器会增加安装、配置、备份和故障排查成本；而SQLite只需要一个数据库文件即可保存结构化数据，Python标准库也直接提供SQLite访问能力。系统可以把数据库文件放在项目数据目录中，备份时复制该文件即可，迁移到其他设备时也更加方便。与此同时，SQLite仍然具备事务能力，可以保证一次业务操作中的多条SQL语句要么一起成功，要么在出错时回滚。

在并发方面，SQLite适合读多写少或轻量写入的场景。考勤系统的写入频率相对有限，主要包括用户注册、人脸提交、管理员审核和签到写入；读取则包括用户中心状态查询、终端人脸库加载和管理员列表查询。本文系统开启外键约束以保证用户、人脸档案和考勤记录之间的关联完整性，并启用WAL日志模式以改善读写并发下的使用体验。系统的数据访问层采用Repository模式封装SQL操作，上层业务不直接操作数据库连接，而是通过Repository完成成员、用户、人脸档案、驳回历史和考勤记录的读写。

在本文系统中，SQLite保存成员名单、用户、人脸档案、驳回记录和考勤记录。成员名单用于约束谁可以注册，用户表用于保存正式账号，人脸档案表用于保存注册照路径、特征向量和审核状态，考勤记录表用于保存签到时间、快照路径和识别置信度。图片文件本身不直接写入数据库，而是保存在本地目录中，数据库只保存路径。这样可以避免数据库文件过大，也便于管理员接口按需读取图片。

\section{Tailscale远程管理网络}

远程管理是树莓派考勤系统的重要辅助能力。考勤终端通常部署在固定位置，管理员不一定能直接接触设备。如果每次审核人脸档案、查询考勤记录或查看设备状态都必须连接树莓派本地显示器和键盘，系统维护成本会明显增加。因此，系统需要一种相对安全、配置简洁的远程访问方式，使Windows管理端能够访问树莓派管理员API。

Tailscale官方文档将其描述为一种基于WireGuard的组网工具，可让设备加入同一个私有网络并进行安全访问\cite{tailscale}。WireGuard是一种现代虚拟专用网络技术，强调配置简洁和加密通信。Tailscale在其基础上提供设备身份、节点发现、密钥管理和私有网络连接能力，使分布在不同网络环境中的设备可以像处在同一内网中一样相互访问。对于没有公网IP或不方便配置端口映射的树莓派，Tailscale可以降低远程访问配置难度。

与直接把树莓派服务暴露到公网相比，Tailscale更适合低成本考勤终端的远程管理场景。公网暴露需要处理端口转发、动态IP、TLS证书、防火墙和攻击面控制等问题；而通过Tailnet访问时，只有加入同一私有网络并具备访问权限的设备才能连接树莓派。本文系统中，树莓派端提供管理员API，Windows管理端通过Tailscale地址访问这些接口。管理员API采用Bearer Token进行鉴权，远程管理端只通过后端服务代理访问树莓派接口和图片资源，浏览器端不直接暴露管理员Token。该设计有利于降低远程管理过程中的凭证泄露和敏感图片直接暴露风险。
